\section{Eleven positive unknowns with denominator reduced absorption}
\label{sec:eleven958}

The new equation has eleven positive unknowns, exact degree $958$, and
$117$ charged binary arithmetic operations. A second equation trades four
degrees for two operations: degree $962$ and $115$ operations. The full
carrier's two-equation system has degrees $(2,956)$ and uses $112$
operations after the standard positive transport shift. These improve
both pinned eleven-positive frontiers, $1154/123$ and $1140/129$.
They do not dominate the twelve-positive degree-$901$ construction.
The argument uses no Sun representation theorem.

Two changes make this possible. First, the multiplier-one norm from
source44 remains sound without the deleted index equation: its main-root
bound needs $d(I\sqrt G-c\sqrt D)>b+c\sqrt D$, not the stronger
$I\sqrt G-c\sqrt D>1$. Second, quotient absorption needs denominator
$Y$, not $E=XY$. The short Pell-index argument recovers the full
$E$-divisibility after the weaker $Y$-divisibility has done its job.

\subsection{Explicit equation}
Take precisely the positive integer unknowns
\[
 J,F,\alpha,z,i,T,s,w,\eta,\zeta,t.
\]
Fix the complete authentic normalized blank-K215 coefficient tuple
$(\beta,K_0,M_C,M_F)$, with $M_F$ already the shifted high mask.
The delivered fixed instance is the materialized Wang943
$121\times8801$ tuple with $\beta=2^{1064921}-1$; its four payloads and
hashes are included. The proof is uniform in any whole authentic tuple.
Set
\begin{gather*}
 q=\beta J+1,\quad s_0=s+3q,\quad X=wq,\quad Y=s_0q^3,\quad E=XY,\\
 k=\eta+\zeta,\quad c=kY+\eta,\quad a=E+Y,\quad H=4a+3,\\
 \Delta=a^2+H=(a+2)^2-1,\quad \ell=EkY,\quad A=1+\ell(\ell+k),\\
 G=1+\Delta c^2,\quad b_0=X+ac,\quad C_w=q-F-\alpha,\\
 R=((q-1)(q-F)+\alpha)(q^2-1)+(M_C+qM_F)J,\\
 N_t=(K_0+w)C_w+q-F-z(q-1),\\
 I=ic^2,\quad D=1+\Delta I^2,\quad Q=\Delta(D-1)=(\Delta I)^2,\\
 d=Q-1,\quad v=cT,\quad b=c+RD,\\
 e=Q(Dv^2+b^2)-1,\quad B_*=2Qvb,\quad K=IT+1,\\
 L_0=1+HK,\quad S=HY,\quad L=YL_0.
\end{gather*}
The computed register $K$ is not the fixed coefficient $K_0$.
For the full carrier put $\delta=(a+2)c-\ell$; for the shorter one put
$\delta=ac-\ell$. In either case define
\[
 W=Y\bigl[b_0+H(t+b+\delta)\bigr]+H(k-R-1)
\]
and the integer polynomial
\begin{equation}\label{eq:new-eleven-norm}
 N=\prod_{\substack{g=\pm\sqrt G\\\tau=\pm\sqrt A\\f=\pm\sqrt D}}
 \left[d(W-Lg-S\tau)^2-S^2e+S^2B_*f\right].
\end{equation}
The formal signs define a conjugate product, not independently assumed
Galois automorphisms. Every other symbol is a polynomial abbreviation,
not another unknown. The proposed single equation is
\begin{equation}\label{eq:new-eleven-single}
 \boxed{\ N_t(1+N)-1=0\ }.
\end{equation}
The two-equation presentation is $N_t=1$, $N=0$.

For an explicit division-free scalar evaluator, set
\begin{align*}
 d_n&=dS^2A,&H_Q&=S^2Q,&H_P&=H_QDv^2,&H_U&=H_Qb^2,\\
 H_e&=H_P+H_U-S^2,&j&=4H_PH_U,\\
 \mathcal T&=2d(W^2+L^2G),&\mathcal K&=[d(W^2-L^2G)]^2,\\
 A_*&=2(d_n+H_e),&C_*&=(d_n-H_e)^2-j,\\
 u&=\mathcal T-A_*,&r&=C_*-\mathcal K,&c_*&=16d_nj,\\
 V_*&=c_*-u(u\mathcal T+2r),&U_*&=r^2-u^2\mathcal K,\\
 N&=\mathcal K V_*^2-U_*(\mathcal T V_*-U_*).
\end{align*}
This is a universal polynomial identity, including degenerate radicands.
The emitted JSON circuits expand every square and count every binary gate.

\subsection{Weighted main-root separation}
Let $x_0=d(W^2+L^2G)$. Since $\mathcal T=2x_0$ and
$\mathcal K\equiv x_0^2\pmod4$, the norm obeys
$N\equiv(U_*-x_0V_*)^2\pmod4$.
Thus a root of \eqref{eq:new-eleven-single} forces $N=0$, $N_t=1$:
the alternative integer unit $1+N=-1$ would make $N=-2$.
The inherited transport interface then gives
\[
 q\ge16,\quad C_w>0,\quad 3q<R<q^4-2.
\]
Our positive shift implies
\[
 X\ge16,\quad Y>3q^4>3R,\quad E>R+1,\quad
 k\ge2,\quad kY<c<2kY,\quad c>96,\quad a+2<2E.
\]
No assertion $k>E$ has been used.

For positive real $g=\sqrt G$, $f=\sqrt D$, $\tau=\sqrt A$, direct
comparison yields
\[
 1<A<\Delta c^2\le\Delta I=\sqrt Q,\qquad
 g>\tau,\quad g<2Ec,\quad f<2EI.
\]
The exact identity $(Ig)^2-(cf)^2=I^2-c^2$ and $I\ge c^2$ give
\[
 Ig-cf>\frac{I}{8Ec}.
\]
Moreover $b+cf<6E^3I^2$ and $d>\tfrac12 E^4I^2$, so
\begin{equation}\label{eq:weighted-gap}
 d(Ig-cf)>\frac{E^3I^3}{16c}
 \ge\frac{E^3I^2c}{16}>6E^3I^2>b+cf.
\end{equation}
This is the new estimate that removes dependence on the missing index
quotient.

Every root of the inner norm has a coherent representation
\[
 x=\tau'+y,\quad (\tau')^2=A,\quad(f')^2=D,\quad
 dy^2=Q(vf'-b)^2-1.
\]
These roots may be complex. With $Z=b+vf>1$,
\[
 |y|<Z+Z/d,\qquad Z/d\le T(b+cf)/d<T(Ig-cf)
\]
by \eqref{eq:weighted-gap}. Hence $|x|<b+(IT+1)g=b+Kg$.
But the negative main sign gives
\[
 \frac{W+Lg}{S}
 =b+t+\delta+\frac{k-R-1}{Y}+\frac{b_0}{H}
       +\left(K+\frac1H\right)g>b+Kg,
\]
since $t+(k-R-1)/Y>0$. This is impossible. Quadratic conjugation
therefore forces the main root $g$ to be a positive integer, and
$x=(W-Lg)/S$ is rational.

\subsection{Integral inner roots}
Set $C(x)=d(x^2+A)-e$. The inner norm is
\[
 F(x)=\bigl[C(x)^2+DB_*^2-4d^2x^2A\bigr]^2
           -4C(x)^2DB_*^2.
\]
If $C(x)=F(x)=0$, substitution gives
$d^2A^2-deA+Q^2v^2b^2D=0$, so $Q\mid A(A-1)$, impossible because
$1<A<\sqrt Q$. Also $F(0)=0$ would imply $Q\mid(A-1)^2$, again
impossible. The rational expressions
\[
 f'=-\frac{C(x)^2+DB_*^2-4d^2x^2A}{2C(x)B_*},\qquad
 \tau'=\frac{C(x)+B_*f'}{2dx},\qquad y=x-\tau'
\]
therefore have justified denominators and satisfy $(f')^2=D$,
$(\tau')^2=A$ and the auxiliary equation. Thus $f',\tau'$ are integers.
The inherited rational Pell saturation lemma,
\[
 m^2V^2-(m^2-1)y^2=1,\quad m,V\in\mathbb Z,\ m\ge2,
 \quad y\in\mathbb Q\quad\Longrightarrow\quad y\in\mathbb Z,
\]
with $m=\Delta I$ and $V=vf'-b$ makes $y$, hence $x$, integral.
A proof of the lemma is included in the companion proof file; it uses
abscissa and ordinate divisibility in the integer-coordinate Pell group.
Reducing $W=Lg+Sx$ modulo $Y$ now yields
\begin{equation}\label{eq:weak-y}
 Y\mid H(k-R-1).
\end{equation}

\subsection{Short indices restore full divisibility}
Write $A_0=a+2$ and let $\chi_{A_0},\psi_{A_0}$ be its Pell sequences.
The main and strong roots give positive indices $p,m$ with
\[
 g=\chi_{A_0}(p),\quad c=\psi_{A_0}(p),\quad
 |f'|=\chi_{A_0}(m),\quad I=\psi_{A_0}(m).
\]
Strong divisibility and $I=ic^2$ give $p\mid m$; composition modulo $c$
then gives $c\mid m/p$, so $m\ge pc>2p$. Also $p<c/2$ and $R<c/2$.
Source38's signed auxiliary-index argument, applied before
\eqref{eq:weak-y}, yields $R\equiv\pm p\pmod c$, hence $p=R$.
For clarity, it classifies the auxiliary Pell solution at $\Delta I$
with odd index $j_0$. Reducing its odd-index polynomial modulo $|f'|$
gives $j_0\equiv\pm p\pmod m$, while reduction modulo $c$ gives
$j_0\equiv\pm R\pmod c$. The full recurrence argument is reproduced
in the companion proof file, including the signed branches.

The square $A$ has fundamental integer-coordinate unit
$(2EY+1,2)$, so for some $n\ge1$,
\[
 |\tau'|=\chi_{2EY+1}(n),\qquad k=2\psi_{2EY+1}(n).
\]
Monotonicity and $2EY+1>A_0$, $c>k$ imply $n<p=R$.
Since $2EY+1\equiv1\pmod E$,
\[
 k\equiv2n\pmod E.
\]
Now $\gcd(Y,H)=\gcd(Y,3)\le3$. By \eqref{eq:weak-y},
\[
 2n\equiv R+1\pmod\mu,\qquad
 \mu=Y/\gcd(Y,H)>R.
\]
But $|2n-R-1|<R$, so $2n=R+1$. It follows that
$E\mid k-R-1$ after all, and
\[
 h=(k-R-1)/E>0.
\]
Dividing the carrier by $Y$ gives
\[
 b_0+H(t+b+\delta+hX)=L_0g+H(\tau'+y).
\]
Thus $g\equiv b_0\pmod H$. The recovered
$\gamma=(g-b_0)/H$ is positive: $n\ge2$ gives $k>X$, hence $c>X$,
and $g>(a+1)c$.
All normalized compiler equations are restored. Absolute values of
$\tau',y$ are harmless; simultaneous sign reversal of $f',T$ is
covered by the inherited signed-quotient interface. This proves soundness
through the calibrated blank-K215 halt equivalence.

\subsection{Completeness and exact degree}
For a canonical compiler witness, $X=2^R$, $Y>2^R$, and $R>3q$
ensure $s_{\rm old}=Y/q^3>3q$. Thus $s=s_{\rm old}-3q>0$.
Choose positive canonical roots and define
\[
 t=\gamma+Kg+\tau+y-b-\delta-hX.
\]
Since $0<\delta<g$, $y\ge b-vf$, and
$\tau>\ell=EkY>k/Y>hX$, one has
\[
 t>\gamma+\tau+T(Ig-cf)-hX>0.
\]
This produces the coherent zero factor and proves completeness.
The transport substitution $z_{\rm old}=z+1$ replaces $N_t=1$ by
$(K_0+w)C_w=F+z(q-1)$ and saves one operation; canonical witnesses
have $z_{\rm old}\ge2$.

Let $N_{12}$ be the earlier multiplier-one norm. The exact localized
identity is
\[
 N=Y^{16}N_{12}\!\left(s_{\rm old}=s+3q,
           t_{\rm old}=t+\frac{k-R-1}{Y}\right).
\]
Each of the eight factors contributes $Y^2$, and the $t_{\rm old}$
degree is at most $16$, so denominators clear. The rational shift has
degree $0$, below the degree-$1$ supplied carrier. It cannot change the
highest form. The invertible highest-linear substitution
$s_{\rm old}=s+3\beta J$ preserves its nonvanishing.
More explicitly, the rational carrier perturbation changes an unscaled
factor only in degrees at most $96$ and $58$, below its smallest
established factor degree $97$.
The earlier exact norm degrees $892$ and $896$ consequently become
$956$ and $960$. The transport factor has exact degree $2$, giving
$958$ and $962$.

The full highest form is $Y_0^{16}$ times
\[
 -2^{14}H_0^{16}d_0^8B^8V^4\ell_0^2k_0Y_0^2
       \left(2t+k_0-\frac{iTc_0}{E_0}\right),
\]
where
\begin{gather*}
 q_0=\beta J,\quad Y_0=(s+3q_0)q_0^3,\quad X_0=wq_0,
 \quad E_0=X_0Y_0,\quad c_0=k_0Y_0,\quad k_0=\eta+\zeta,\\
 I_0=ic_0^2,\quad H_0=4E_0,\quad \ell_0=E_0c_0,
 \quad d_0=E_0^4I_0^2,\\
 B=q_0^3(q_0-F)E_0^2I_0^2,\qquad V=I_0T\ell_0.
\end{gather*}
Its $t$ coefficient is nonzero for every authentic $\beta>0$;
the displayed denominator cancels. Multiplying by
$w(q_0-F-\alpha)-zq_0$ gives the single equation's highest form.
The companion proof gives the short carrier's form as well.

\subsection{Verification and logical boundary}
The included builder audits all gates and coefficient hashes. The
independent checker passes 24 exact and 60 modular formal conjugate
products per carrier, three complete finite-field polynomial rays with
both illustrative and actual fixed coefficients, all stated leading
forms, the transport shift, and 288 pre-index bound tests. The tests
include $k=2$ and arbitrarily large $X$, where the old unweighted gap
shortcut is unavailable. Both normal Python and optimized Python pass.
Finite checks validate the implementation, not the unbounded compiler.

The baseline is ProveIt commit
\texttt{124c07c0b77a}, specifically source38's
index reconstruction and source44's multiplier-one norm. Complete source
pins, copied proof dependencies, payloads, circuits, and the detailed
mathematical proof accompany this section. The resulting theorem is
that ZFC proves the absence of positive solutions equivalent to
$\operatorname{Con}(\mathrm{ZFC})$. No consistency decision, global
minimality theorem, independent human review, or new kernel-checked
end-to-end certificate is asserted.
